\chapter{Technical Background}
\label{ch:pfe-background}

\section{Introduction}

Recovering a Software Bill of Materials from a compiled \Android application
connects several technical domains. The task begins with software
supply-chain security, because an application inherits code and risk from
external components. It then requires a representation for those components,
a method for relating them to vulnerability information, and an understanding
of how the \Android build process transforms source dependencies into an
\APK.

This chapter develops that chain in the same order. It first introduces
software supply chains, \SBOM concepts, package identifiers, and vulnerability
terminology. It then explains the \Android ecosystem, the internal structure
of an \APK, the D8 and \Reight compilation stages, and the evidence available
to static reverse engineering. The final sections distinguish library
identity from exact-version identification and define the evidence,
provenance, confidence, and uncertainty concepts used throughout this thesis.

The purpose is not to reproduce the implementation described in
Chapters~\ref{ch:pfe-design} and~\ref{ch:pfe-evaluation}. It is to establish
why a compiled application cannot be treated as a dependency manifest and why
trustworthy recovery must combine several imperfect evidence sources.

\section{Software Supply-Chain Security}

\subsection{Software Supply Chains}

A software supply chain comprises the people, processes, tools, services, and
artifacts involved in producing and delivering software. It includes source
repositories, build systems, package registries, external libraries,
continuous-integration services, signing infrastructure, distribution
channels, and the final deployed artifact. NIST treats cybersecurity
supply-chain risk management as a life-cycle activity because a weakness or
compromise may enter through any of these elements, including products and
services acquired from external suppliers \cite{nist2022csc}.

Third-party dependencies are therefore part of the application's security
boundary. They reduce development effort and provide mature functionality,
but they also introduce inherited code, licenses, maintenance schedules, and
vulnerabilities. A single dependency may itself depend on other components,
forming a transitive graph rather than a flat list.

For an open-source project, a build file and resolved dependency graph can
describe this supply-chain view before compilation. A released binary offers
a different view: it contains the code and resources that survived resolution,
merging, shrinking, and packaging. These two views overlap, but they are not
identical. A declared dependency may contribute no code to the final artifact,
while copied, shaded, vendored, generated, or dynamically obtained code may
not appear as a simple direct dependency row.

\subsection{Supply-Chain Attacks}

Supply-chain attacks target trust relationships instead of attacking only the
final application. Relevant attack paths include:

\begin{itemize}
  \item compromising a source repository or developer account;
  \item publishing a malicious or dependency-confusion package;
  \item modifying an artifact in a registry or distribution channel;
  \item compromising a build service, compiler, plug-in, or signing key;
  \item exploiting a known vulnerability in a reused component; and
  \item inserting or downloading code outside the expected build graph.
\end{itemize}

These cases have different causes, but they share an attribution problem:
security teams need to know which component is present, which version or
variant is involved, where it originated, and which evidence supports the
claim. A component name without a version is often insufficient for
vulnerability assessment, while a version without provenance may be
impossible to audit.

\subsection{Transparency, Traceability, and Provenance}

\emph{Transparency} makes the composition and production context of software
visible. \emph{Traceability} connects a delivered artifact to upstream
components and transformations. \emph{Provenance} records where an artifact or
claim came from and how it was produced. These properties support incident
response, patch prioritization, supplier assessment, and reproducible
analysis.

In this thesis, provenance is also applied to analysis evidence. A recovered
component claim should preserve at least:

\begin{itemize}
  \item the analyzed \APK and its cryptographic hash;
  \item the observed file, \DEX, class, method, literal, or structural signal;
  \item the reference artifact and version used for comparison;
  \item the matching or validation procedure;
  \item the confidence and competing alternatives; and
  \item the evidence class, such as controlled, source-validated, partial, or
  experiment-only.
\end{itemize}

This is narrower than full build provenance, because a binary-only analyzer
cannot recreate every upstream build event. It is nevertheless essential for
making the resulting \SBOM auditable.

\section{Software Bills of Materials}

\subsection{Definition and Purpose}

An \SBOM is a structured inventory of software components and their
relationships. The NTIA describes it as a nested inventory for software,
analogous to a list of ingredients, and emphasizes its role in software
transparency and vulnerability management \cite{ntiaSbom}. An \SBOM can be
created during a trusted build, supplied by a vendor, generated from source
metadata, or reconstructed from a compiled artifact.

A useful component record commonly contains:

\begin{itemize}
  \item a component type, supplier, name, namespace, and version;
  \item one or more persistent package identifiers;
  \item cryptographic hashes;
  \item licenses and copyright information;
  \item dependency or containment relationships;
  \item external references and vulnerability links; and
  \item evidence or properties that explain how the record was obtained.
\end{itemize}

Not every field is available in every production context. A source-generated
\SBOM may know the declared dependency graph and licenses but not which methods
survived optimization. A binary-recovered \SBOM observes the shipped artifact
but may be uncertain about coordinate, version, scope, or transitive
relationships. The generation method is therefore part of the meaning of the
document.

\subsection{\SPDX and \CycloneDX}

\SPDX and \CycloneDX are widely used machine-readable standards for software
composition data. \SPDX originated around licensing and package information
and now supports several software supply-chain profiles. \CycloneDX is an
OWASP standard designed for bill-of-materials and risk-analysis use cases,
including components, services, dependencies, vulnerabilities, compositions,
and evidence \cite{spdxSpecification,cyclonedxSpecification}.

Table~\ref{tab:pfe-background-sbom-standards} compares the dimensions relevant
to this project. The comparison does not imply that one standard is universally
superior; the choice depends on the information available and the downstream
workflow.

\begin{table}[H]
\centering
\caption{SBOM-standard dimensions relevant to compiled APK analysis.}
\label{tab:pfe-background-sbom-standards}
\begin{tabularx}{\textwidth}{L{0.20\textwidth}XX}
\toprule
\textbf{Dimension} & \textbf{\SPDX} & \textbf{\CycloneDX}\\
\midrule
Primary emphasis & Software-package, file, license, security, build, and
provenance information through defined profiles & Component and service
inventory, dependency relationships, security analysis, vulnerabilities, and
operational risk\\
Representations & JSON, JSON-LD, RDF, YAML, XML, and tag/value depending on
the specification version & JSON, XML, and Protocol Buffers depending on the
specification version\\
Identity & Package metadata, external references, checksums, and package
identifiers & Component metadata, Package URL, CPE, hashes, properties, and
external references\\
Relationships & Rich typed relationships among elements & Dependency graph,
composition, formulation, and service relationships\\
Evidence use & Supports provenance and declared or concluded information &
Supports component evidence, identity evidence, occurrences, methods, and
confidence-oriented properties\\
Project choice & Suitable alternative or exchange format & Selected output
format because it maps directly to component, PURL, evidence, dependency, and
vulnerability workflows\\
\bottomrule
\end{tabularx}
\end{table}

The MVP emits \CycloneDX JSON. This is an engineering choice, not a claim that
\SPDX is unsuitable. A later exporter could map sufficiently supported
component data to either format.

\subsection{Package URL}

Package URL is a standardized package identifier whose general form is:

\begin{center}
\small\path{pkg:type/namespace/name@version?qualifiers#subpath}
\end{center}

It identifies software according to a package ecosystem rather than by a
free-text name \cite{packageUrlSpecification}. For a Maven component, a PURL
can represent the group, artifact, and version, for example:

\begin{center}
\texttt{pkg:maven/androidx.fragment/fragment@1.6.1}
\end{center}

This normalization is important because names such as ``fragment'' or
``core'' are not globally unique. The ecosystem and namespace reduce
ambiguity and give vulnerability services a stable query key. However, a
syntactically valid PURL does not prove that the identified component exists
in an \APK. Its trustworthiness remains bounded by the evidence used to infer
the coordinate and version.

\subsection{Use Cases and Limitations}

\SBOM{}s support component inventory, vulnerability triage, license review,
supplier communication, incident response, asset correlation, and policy
checking. During a newly disclosed vulnerability, an organization can search
its inventories for affected packages instead of manually examining every
application.

An \SBOM is not, by itself, proof that software is secure. It may be incomplete,
outdated, generated with weak identifiers, or disconnected from the deployed
artifact. It does not establish exploitability, configuration, reachability,
or maliciousness. For binary reconstruction, the limitations are stronger:
shrinking can remove dependency evidence; optimization can mix boundaries;
native code may lack ecosystem metadata; and runtime-loaded components may be
absent from the package. A responsible recovered \SBOM must expose these
limitations instead of presenting inferred rows as original build facts.

\section{Vulnerability Identification and Management}

\subsection{Identifiers, Weaknesses, Scores, and Catalogs}

Vulnerability management uses several related but non-interchangeable
concepts. Table~\ref{tab:pfe-background-vulnerability-vocabulary} summarizes
the role of each one.

\begin{table}[H]
\centering
\caption{Vulnerability concepts used by the APK-to-SBOM workflow.}
\label{tab:pfe-background-vulnerability-vocabulary}
\begin{tabularx}{\textwidth}{L{0.13\textwidth}L{0.25\textwidth}X}
\toprule
\textbf{Term} & \textbf{Role} & \textbf{Interpretation}\\
\midrule
CVE & Public vulnerability identifier & Names a publicly disclosed
vulnerability record so tools and organizations can refer to the same issue
\cite{cveOverview}.\\
CWE & Weakness taxonomy & Classifies recurring design and implementation
weaknesses that may lead to vulnerabilities; it is not a component-version
advisory list \cite{cweAbout}.\\
NVD & Enriched vulnerability database & Adds analysis and structured data to
CVE records, including affected-product and severity information
\cite{nvdGeneral}.\\
CVSS & Technical severity framework & Describes intrinsic and environmental
severity characteristics; it does not predict whether exploitation will occur
\cite{firstCvss}.\\
EPSS & Exploitation-probability model & Estimates the probability that a
published CVE will be exploited in the wild during the next 30 days
\cite{firstEpss}.\\
KEV & Known-exploitation catalog & Lists vulnerabilities for which CISA has
evidence of exploitation and provides remediation-oriented operational
priority \cite{cisaKev}.\\
OSV & Ecosystem-aware advisory schema and service & Represents affected
package ecosystems and version ranges and supports package/version queries
\cite{osvSchema,osvApi}.\\
\bottomrule
\end{tabularx}
\end{table}

These sources answer different questions. CVSS concerns severity, EPSS concerns
forecast exploitation probability, and KEV records known exploitation. None
of them repairs an incorrect package identity. The first requirement is still
to determine whether an affected component and version are present.

\subsection{Why Exact Identity and Version Matter}

Most vulnerability advisories describe an affected package and one or more
version ranges. Mapping the advisory to an application therefore depends on
the chain:

\begin{equation}
\label{eq:pfe-background-vulnerability-chain}
\text{binary evidence}
\rightarrow
\text{package coordinate}
\rightarrow
\text{version}
\rightarrow
\text{affected range}
\rightarrow
\text{vulnerability context}.
\end{equation}

An error early in the chain propagates. A wrong coordinate can produce
unrelated advisories; a wrong version can create a false positive or hide a
real exposure. Adjacent library versions are especially difficult because
they may differ by only a few methods, and optimization can remove precisely
those differences.

This thesis consequently distinguishes three forms of vulnerability output:

\begin{itemize}
  \item \textbf{confirmed mapping}: the component and exact version have
  sufficiently strong evidence, and the version falls within an advisory's
  affected range;
  \item \textbf{potential mapping}: the component is plausible, but version,
  variant, reachability, or affected-range evidence is incomplete; and
  \item \textbf{report-only enrichment}: external risk information is shown
  for analyst review but is not allowed to change the authoritative \SBOM or
  source-validation metrics.
\end{itemize}

Even a confirmed package/version mapping is not a proof of runtime
exploitability. Configuration, platform behavior, call reachability, and
mitigations still require separate assessment.

\section{Android Application Ecosystem}

\subsection{Application Architecture and Development}

\Android applications are commonly written in Kotlin or Java, with optional
C or C++ native code, and are executed within the \Android application
framework and runtime. Application components such as activities, services,
broadcast receivers, and content providers are declared through the
application manifest and interact with platform services through framework
APIs \cite{androidAppFundamentals,androidManifestOverview}.

The standard build system uses Gradle with the Android Gradle Plugin.
Application code, generated code, resources, local modules, and external
dependencies are processed into an installable artifact
\cite{androidBuildOverview}. Distribution may use an Android App Bundle,
from which an app store generates device-specific APKs, while direct analysis
usually receives one universal or split \APK.

\subsection{JAR, AAR, and Maven Dependencies}

Java libraries are commonly distributed as JAR files containing compiled JVM
classes and metadata. Android libraries are commonly distributed as AAR
archives, which may contain a JAR together with Android resources, a manifest,
native libraries, consumer rules, and other Android-specific content
\cite{androidLibrary}. Maven coordinates conventionally identify an artifact
as \texttt{group:artifact:version}.

Dependency resolution is variant-aware. Gradle selects compatible variants
according to attributes such as platform, usage, build type, flavor, and
target environment \cite{gradleVariantResolution}. This matters for reverse
identification because a public dependency coordinate can resolve to
different physical artifacts:

\begin{itemize}
  \item a base Java or Android artifact;
  \item a Kotlin extension artifact, often identified by a \texttt{-ktx}
  suffix;
  \item a JVM-, Android-, or other platform-specific target;
  \item a metadata or bill-of-materials artifact with no runtime bytecode; or
  \item a shaded or repackaged artifact whose internal classes no longer
  preserve the public coordinate name.
\end{itemize}

The analyzer observes the selected and transformed binary content, not the
declarative choice that produced it. Public-coordinate, target-variant, and
packaging differences must therefore be considered when building reference
profiles.

\section{APK Structure}

\subsection{The APK as a Container}

An \APK is a ZIP-format package containing executable code, resources,
configuration, signatures, and optional native binaries. Android Studio's APK
Analyzer exposes these archive entries and reports DEX, resource, and file
composition information \cite{androidApkAnalyzer}. Figure
\ref{fig:pfe-background-apk-structure} presents the main evidence-bearing
regions.

\begin{figure}[H]
\centering
\begin{tikzpicture}[
  node distance=5mm,
  box/.style={draw, rounded corners, align=left, text width=0.72\textwidth,
    minimum height=8mm, fill=blue!4},
  title/.style={draw, rounded corners, align=center, text width=0.72\textwidth,
    minimum height=9mm, fill=schoolblue!15, font=\bfseries}
]
\node[title] (apk) {Android application package};
\node[box, below=of apk] (manifest)
  {\texttt{AndroidManifest.xml}: package, components, permissions, SDK and
  application metadata};
\node[box, below=of manifest] (dex)
  {\texttt{classes.dex}, \texttt{classes2.dex}, \ldots: application and
  embedded JVM-language bytecode};
\node[box, below=of dex] (resources)
  {\texttt{resources.arsc}, \texttt{res/}, and \texttt{assets/}: compiled and
  raw resources, strings, configuration, and bundled data};
\node[box, below=of resources] (native)
  {\texttt{lib/<ABI>/*.so}: architecture-specific native shared libraries};
\node[box, below=of native] (meta)
  {\texttt{META-INF/} and signing blocks: integrity and signer evidence};
\node[box, below=of meta] (other)
  {Other files: configuration, models, archives, secondary payloads, or
  nonstandard code-loading material};
\end{tikzpicture}
\caption{Principal APK regions and the evidence they can expose.}
\label{fig:pfe-background-apk-structure}
\end{figure}

\subsection{Manifest, Resources, and Assets}

The manifest declares the application package, components, permissions,
features, SDK constraints, intent filters, and other installation or runtime
metadata \cite{androidManifestOverview}. In a packaged \APK it is stored in a
compiled binary XML representation. Manifest metadata can support application
identification and contextual analysis, but it does not normally enumerate
all embedded Maven dependencies.

Compiled resources are indexed through \texttt{resources.arsc} and the
\texttt{res/} hierarchy. The \texttt{assets/} directory may contain arbitrary
application data. Resource names, strings, notices, and configuration files
can occasionally reveal component evidence; however, these signals may be
renamed, removed, shared among products, or unrelated to executable code.

\subsection{DEX and Multidex}

\DEX is the executable bytecode format used by the \Android runtime. A DEX
file contains type, string, field, method, class, annotation, and instruction
data. Application and dependency classes are typically merged into one or
more DEX files, so the original JAR or AAR boundaries are no longer present.

The DEX format imposes reference-index limits. Applications exceeding the
single-DEX method-reference limit use multidex packaging, commonly producing
\texttt{classes.dex}, \texttt{classes2.dex}, and additional numbered files
\cite{androidMultidex}. A component detector must analyze all relevant DEX
entries. Examining only the primary DEX can silently miss dependencies.

\subsection{Signing and Native Libraries}

\Android requires APKs to be digitally signed. The signing identity supports
update trust and package integrity; modern signing schemes can protect the
whole artifact rather than only individual ZIP entries
\cite{androidSigning}. A signer does not identify every embedded dependency,
but signing information is useful provenance for the application artifact.

Native libraries are stored under ABI-specific directories such as
\texttt{lib/arm64-v8a/} or \texttt{lib/x86\_64/}. An ABI defines the binary
interface expected by a CPU architecture and platform
\cite{androidAbi}. Native file names and hashes are observable, but general
package and version attribution from stripped shared objects is a separate
binary-analysis problem. This thesis therefore distinguishes native file
evidence from Maven component identification.

\subsection{Nonstandard and Dynamically Loaded Code}

Not all executable behavior must appear in conventional
\texttt{classes*.dex} entries. Applications may unpack encrypted payloads,
load additional DEX or native code, use custom class loaders, or download code
after installation. Android security guidance warns that dynamic code loading
increases risk, especially when remote sources or weak integrity controls are
used \cite{androidDynamicCode}.

A static \APK analyzer can report suspicious archives, nonstandard DEX paths,
native-heavy packages, and other observable indicators. It cannot truthfully
inventory code that is absent, encrypted beyond available analysis, or fetched
only at runtime. These are scope limitations, not ordinary negative detection
results.

\section{Android Compilation and Optimization}

\subsection{From Source Dependencies to APK}

Figure~\ref{fig:pfe-background-build-pipeline} summarizes the build path that
turns source and external artifacts into packaged evidence.

\begin{figure}[H]
\centering
\begin{tikzpicture}[
  node distance=7mm and 7mm,
  stage/.style={draw, rounded corners, align=center, minimum height=10mm,
    text width=0.19\textwidth, fill=blue!5},
  transform/.style={draw, rounded corners, align=center, minimum height=10mm,
    text width=0.16\textwidth, fill=warningorange!12},
  arrow/.style={-{Latex[length=2.5mm]}, thick}
]
\node[stage] (source) {Kotlin/Java source\\generated source};
\node[stage, right=of source] (deps) {Resolved JAR/AAR\\dependency variants};
\node[transform, right=of deps] (class) {JVM bytecode\\and resources};
\node[transform, below=of class] (d8) {D8 or \Reight\\DEX transformation};
\node[stage, left=of d8] (package) {DEX, manifest,\\resources, native files};
\node[stage, left=of package] (signed) {Packaged and\\signed APK};

\draw[arrow] (source) -- (deps);
\draw[arrow] (deps) -- (class);
\draw[arrow] (class) -- (d8);
\draw[arrow] (d8) -- (package);
\draw[arrow] (package) -- (signed);
\end{tikzpicture}
\caption{Simplified Android build path from source and resolved dependencies
to a signed APK.}
\label{fig:pfe-background-build-pipeline}
\end{figure}

Kotlin and Java sources are compiled to JVM bytecode. Resources and manifests
are processed separately, while external artifacts are resolved according to
the build's dependency and variant rules. D8 transforms Java bytecode into DEX
bytecode and can receive compiled classes, JARs, APKs, and ZIP archives as
inputs \cite{androidD8}. Packaging then combines code, resources, native
libraries, and metadata before signing.

This process destroys a convenient assumption for component recovery: the
output is not a collection of untouched dependency archives. Classes from the
application and many libraries coexist in merged DEX files, while AAR
resources are merged into the application resource table.

\subsection{R8 and ProGuard-Compatible Rules}

\Reight combines shrinking, optimization, obfuscation, and bytecode
transformation in the Android release build process
\cite{androidShrinkCode}. It consumes keep rules, including rules written in
the ProGuard configuration language, but \Reight is more than a simple
identifier renamer.

Its transformations include:

\begin{itemize}
  \item \textbf{renaming}: replacing class, method, and field names with short
  identifiers;
  \item \textbf{shrinking}: removing unreachable classes, methods, fields, and
  resources;
  \item \textbf{optimization}: rewriting instructions and simplifying control
  or data flow;
  \item \textbf{inlining}: moving a method body into its callers and removing
  the original boundary;
  \item \textbf{class merging or repackaging}: changing structural and package
  organization; and
  \item \textbf{constant and resource propagation}: replacing or eliminating
  values that would otherwise provide literals or references.
\end{itemize}

Figure~\ref{fig:pfe-background-r8-evidence} relates these transformations to
the evidence used by a component detector.

\begin{figure}[H]
\centering
\begin{tikzpicture}[
  node distance=8mm and 11mm,
  before/.style={draw, rounded corners, align=center, text width=0.25\textwidth,
    minimum height=11mm, fill=evidencegreen!10},
  action/.style={draw, rounded corners, align=center, text width=0.25\textwidth,
    minimum height=11mm, fill=warningorange!15},
  after/.style={draw, rounded corners, align=center, text width=0.25\textwidth,
    minimum height=11mm, fill=red!7},
  arrow/.style={-{Latex[length=2.5mm]}, thick}
]
\node[before] (names) {Readable package and\\class descriptors};
\node[action, right=of names] (rename) {Renaming and\\repackaging};
\node[after, right=of rename] (namesloss) {Name overlap weakens};
\node[before, below=of names] (methods) {Library methods and\\class boundaries};
\node[action, right=of methods] (shrink) {Shrinking, inlining,\\and class merging};
\node[after, right=of shrink] (methodloss) {Coverage and graph\\boundaries change};
\node[before, below=of methods] (ops) {Opcodes, literals, calls,\\and dependencies};
\node[action, right=of ops] (rewrite) {Optimization and\\constant propagation};
\node[after, right=of rewrite] (oploss) {Some low-level signals\\survive; others drift};

\draw[arrow] (names) -- (rename) -- (namesloss);
\draw[arrow] (methods) -- (shrink) -- (methodloss);
\draw[arrow] (ops) -- (rewrite) -- (oploss);
\end{tikzpicture}
\caption{How R8 transformations weaken different component-identification
signals.}
\label{fig:pfe-background-r8-evidence}
\end{figure}

Renaming mainly attacks textual identity, while shrinking and inlining alter
the amount and organization of surviving code. Two neighboring versions may
become observationally identical if \Reight removes their version-distinguishing
methods. Conversely, a heavily transformed component may retain distinctive
opcode, literal, call, or dependency patterns even after names disappear.
This motivates multi-feature matching and explicit abstention.

\section{Reverse Engineering and Static Analysis}

\subsection{Static and Dynamic Analysis}

Static analysis examines an application without executing it. For Android,
this may include archive inspection, manifest decoding, resource analysis,
DEX disassembly, control-flow construction, call analysis, and native binary
inspection. OWASP describes static analysis as examining an application's
code and resources to identify properties without runtime execution
\cite{owaspStaticAnalysis}.

Dynamic analysis observes behavior while the application runs, often through
instrumentation, debugging, system-call observation, network interception, or
runtime hooking. It can reveal decrypted or downloaded code and actual
execution paths, but its coverage depends on devices, inputs, permissions, and
the explored behavior. Reverse engineering commonly combines both approaches
when the objective requires it \cite{owaspReverseEngineering}.

The MVP in this thesis is intentionally static. Static analysis is
reproducible, suitable for offline batch processing, and directly connected
to the shipped \APK. The trade-off is that purely runtime or protected code
may remain invisible.

\subsection{Recoverable DEX Evidence}

A DEX parser can recover several evidence families:

\begin{itemize}
  \item package and class descriptors;
  \item fields and method signatures;
  \item instruction opcodes and opcode-frequency summaries;
  \item string literals and constants;
  \item method invocations and framework API references;
  \item inheritance, interface, field-type, and call dependencies; and
  \item class-, method-, or basic-block-level graph structure.
\end{itemize}

These signals have different robustness. Class names can be highly
discriminative in unobfuscated applications but fragile under renaming.
Opcode distributions are less dependent on names but may be too coarse for
nearby versions. Literals can distinguish releases, yet optimization may
remove them. Dependency graphs preserve structural context, but inlining and
class merging can change their topology. No feature family should therefore
be interpreted independently of its failure modes.

\subsection{Reverse-Engineering Challenges}

Practical analysis must handle multidex applications, malformed entries,
unusual archive paths, generated code, duplicate classes, library repackaging,
Kotlin compiler artifacts, reflection, resource-only dependencies, native
code, and protection systems. It must also distinguish a parse failure from an
empty result. If one DEX file cannot be analyzed, silently treating it as
containing no components would convert missing evidence into false negative
evidence.

Legal and ethical boundaries also matter. Reverse engineering should be
performed on authorized artifacts for research, interoperability, audit, or
security purposes consistent with applicable law and policy. This project
uses static analysis to identify components and evidence; it does not attempt
to bypass access controls or modify third-party applications.

\section{Dependency and Component Identification}

\subsection{Library Identity versus Version Identity}

Third-party library identification asks whether code from a known component is
present. Version identification asks which release or release range best
explains the observed evidence. The second task is strictly harder. Many
versions share most packages, classes, methods, and strings, while a small
patch may change only a few instructions.

This distinction is central to \SBOM and vulnerability accuracy. A detector
may have strong evidence for \texttt{androidx.fragment:fragment} while lacking
enough surviving evidence to separate versions \texttt{1.6.1} and
\texttt{1.6.2}. Reporting the library identity can be justified while
asserting an exact version cannot.

\subsection{Reference Profiles}

Reference-based detection begins with known artifacts. For each coordinate
and version, the system extracts a profile containing comparable evidence,
such as descriptors, methods, opcodes, literals, API references, and
dependency structure. The same extraction process is applied to the target
\APK. Comparing DEX-like representations is preferable to comparing source
code directly with optimized DEX because it reduces representation mismatch.

Reference coverage creates an important boundary. If the correct component or
version is absent from the reference database, the matcher cannot select it.
A high similarity to the nearest available version is not proof that the
available version is correct. Reference absence must be represented as a
coverage limitation rather than hidden inside the matcher score.

\subsection{Candidate Generation, Similarity, and Confidence}

A scalable matcher commonly uses two stages. Candidate generation removes
obviously unrelated references through inexpensive evidence such as package
or class overlap. Detailed scoring then combines several feature similarities
for the remaining candidates. This coarse-to-fine pattern reduces cost, but an
overly strict shortlist can discard the correct component after obfuscation.

A weighted score can be expressed abstractly as:

\begin{equation}
\label{eq:pfe-background-weighted-score}
S(c,v \mid A) = \sum_{i=1}^{n} w_i s_i(A,R_{c,v}),
\end{equation}

where \(A\) is the analyzed application, \(R_{c,v}\) is the reference profile
for coordinate \(c\) and version \(v\), \(s_i\) is a feature-specific
similarity, and \(w_i\) is its configured weight. The score ranks evidence; it
is not automatically a calibrated probability that the component is present.

Confidence also depends on separation from alternatives, feature coverage,
evidence quality, reference completeness, and known transformations.
Table~\ref{tab:pfe-background-identity-version} illustrates the required
decision boundary.

\begin{table}[H]
\centering
\caption{Example distinction between component and exact-version evidence.}
\label{tab:pfe-background-identity-version}
\begin{tabularx}{\textwidth}{L{0.23\textwidth}L{0.20\textwidth}XX}
\toprule
\textbf{Observed situation} & \textbf{Identity claim} &
\textbf{Version claim} & \textbf{Permitted output}\\
\midrule
Distinctive classes, methods, literals, and structure match one reference &
Strong & Strong and separated & Emit coordinate and exact version with
supporting evidence\\
Library family is clear, but two adjacent versions have nearly equal evidence
& Strong & Ambiguous & Emit identity only or preserve both version candidates;
do not invent a tiebreak\\
Names are removed, while partial structural and opcode evidence remains &
Candidate & Weak & Report as uncertain or abstain from authoritative SBOM
emission\\
No suitable reference exists & Unknown & Unsupported & Record a reference
coverage gap, not component absence\\
Only a native file hash is visible & File observation & Unsupported Maven
version & Emit a native file component if useful, without claiming a Maven
coordinate\\
\bottomrule
\end{tabularx}
\end{table}

\section{Evidence-Based SBOM Recovery}

\subsection{Evidence Concepts}

Evidence-based recovery treats every emitted component as a claim that must be
connected to observable support. Five concepts organize that support:

\begin{itemize}
  \item \textbf{evidence source}: the artifact or observation from which a
  signal was obtained;
  \item \textbf{provenance}: the chain linking the signal to the analyzed
  \APK, reference artifact, tool, configuration, and procedure;
  \item \textbf{strength}: how specifically and reliably the signal supports
  an identity or version;
  \item \textbf{conflict}: disagreement among sources or similarly supported
  candidates; and
  \item \textbf{uncertainty}: the bounded lack of knowledge that remains after
  analysis.
\end{itemize}

Evidence strength is contextual. A readable package name may be strong in a
clear application and weak after repackaging. A shared literal may be weak
alone but useful when combined with version-unique methods and dependency
structure. Evidence must therefore be interpreted as a set rather than a
collection of universally reliable indicators.

\subsection{Evidence Sources and Failure Modes}

Table~\ref{tab:pfe-background-evidence-sources} summarizes the main evidence
sources available to a static APK analyzer.

\begin{table}[H]
\centering
\caption{Static evidence sources for APK-derived SBOM recovery.}
\label{tab:pfe-background-evidence-sources}
\begin{tabularx}{\textwidth}{L{0.19\textwidth}L{0.29\textwidth}X}
\toprule
\textbf{Source} & \textbf{Useful evidence} & \textbf{Principal limitation}\\
\midrule
APK container & File paths, sizes, hashes, DEX inventory, native files,
signing and suspicious payload indicators & Archive presence does not provide
package identity for every file\\
Manifest & Application identity, components, permissions, SDK and feature
context & Normally does not enumerate embedded Maven dependencies\\
DEX descriptors & Packages, classes, fields, methods, types & Renaming,
repackaging, merging, generated code, and multidex omissions\\
DEX instructions & Opcodes, calls, constants, strings, API references &
Optimization, inlining, dead-code removal, and common compiler patterns\\
Dependency structure & Class, method, inheritance, field-type, and call
relationships & Graph topology changes under shrinking, merging, and inlining\\
Resources and assets & Resource names, notices, metadata, bundled files and
strings & Shared or unrelated data; removal and renaming; no guaranteed
version semantics\\
Native libraries & ABI, filename, size, hash, imports, and binary structure &
Stripping, static linking, absent ecosystem identifiers, and separate native
version-analysis requirements\\
Reference artifacts & Known coordinate/version profiles and version-unique
evidence & Catalog incompleteness, packaging errors, variant mismatch, and
unavailable historical releases\\
Source validation & Resolved Gradle dependency graph for an open-source app &
Represents the source/build view and may differ from code surviving in the
distributed APK\\
\bottomrule
\end{tabularx}
\end{table}

\subsection{Fusion, Conflict, and Abstention}

Evidence fusion combines signals that fail differently. Names can establish
clear identities, structural evidence can survive some renaming, literals can
separate particular versions, and source validation can verify emitted rows
for open-source applications. Agreement among independent signals supports a
stronger claim than repetition of one signal family.

Conflict should not be hidden by a forced maximum. If one candidate wins by
class overlap while another wins by version-unique literals, the report should
preserve the disagreement and identify which signal may have been affected by
optimization or catalog error. Similarly, a narrow score margin between
adjacent versions is evidence of ambiguity, not a reason to select the newer
release.

Abstention is therefore a valid system outcome. The analyzer may:

\begin{itemize}
  \item emit a coordinate and exact version;
  \item emit a coordinate while withholding exact version authority;
  \item expose ranked candidates for analyst review;
  \item emit only a hash-based native file observation;
  \item record a reference or parse coverage gap; or
  \item classify the case as outside the supported static-analysis boundary.
\end{itemize}

This hierarchy prevents uncertainty from being converted into false precision.
It also protects vulnerability analysis: a potential candidate may support a
warning or investigation path, but it must not silently become a confirmed
\CVE finding.

\section{Conclusion}

Software supply-chain security requires knowledge of the components delivered
inside an application, yet the compiled \APK does not retain a complete
dependency manifest. JAR and AAR boundaries are merged, DEX files combine
application and library code, resources are processed, and \Reight may rename,
remove, inline, or reorganize the very evidence needed for attribution.

\SPDX, \CycloneDX, and Package URL provide standard structures for describing
components, but a standard representation cannot compensate for weak
identity or version evidence. Vulnerability databases likewise depend on
correct package/version mapping and must remain subordinate to the certainty
of the recovered component claim.

Consequently, no single APK signal is sufficient for trustworthy recovery.
Names, methods, opcodes, literals, dependency structure, native hashes,
reference profiles, and source validation each contribute useful but bounded
evidence. A defensible system must fuse these signals, retain provenance,
expose conflicts, and abstain when exact identification is not supported. The
next chapter examines whether machine-learning and deep-learning methods can
assist this process while remaining outside the authoritative production
boundary.
